% Thesis-native adaptation; the original report is in source_as_supplied/.
% References are collected in XXX_Bibliography/Appendix_SpalartAllmaras.bib.
\chapter[Density treatment in the Spalart--Allmaras model]{Consistency of the density treatment in the compressible Spalart--Allmaras model of OpenFOAM v2412}
\label{app:spalart-allmaras-density}

\begingroup
\definecolor{SAHypBlue}{HTML}{17324D}
\definecolor{SAHypSteel}{HTML}{456C86}
\definecolor{SAHypLightBlue}{HTML}{EAF2F7}
\definecolor{SAHypLightGreen}{HTML}{EAF6EF}
\definecolor{SAHypLightRed}{HTML}{FBEDEC}
\definecolor{SAHypLightGold}{HTML}{FFF6E5}
\definecolor{SAHypGray}{HTML}{F5F7F9}
\definecolor{SAHypRule}{HTML}{D8E0E6}

\setlength{\parindent}{0pt}
\setlength{\parskip}{0.65em}
\setlength{\emergencystretch}{2em}
\renewcommand{\arraystretch}{1.22}
\setlist[itemize]{leftmargin=*,topsep=0.35em,itemsep=0.2em}
\setlist[enumerate]{leftmargin=*,topsep=0.35em,itemsep=0.3em}
\makeatletter
\setlength{\@fptop}{0pt}
\makeatother

\newcommand{\SADt}{\frac{\mathrm D}{\mathrm Dt}}
\newcommand{\SAgrad}{\boldsymbol{\nabla}}
\newcommand{\SAdiv}{\boldsymbol{\nabla}\!\cdot}
\newcommand{\SAcode}[1]{\texttt{#1}}
\newcommand{\SAOF}{OpenFOAM}

\newcommand{\SAnote}[2]{%
  \vspace{0.3em}%
  \noindent\colorbox{#1}{%
    \begin{minipage}{\dimexpr\linewidth-2\fboxsep\relax}
      #2
    \end{minipage}}%
  \vspace{0.3em}%
}

\newcommand{\SAeqbox}[2]{%
  \vspace{0.1em}%
  \noindent\fcolorbox{SAHypRule}{#1}{%
    \begin{minipage}{\dimexpr\linewidth-2\fboxsep-2\fboxrule\relax}
      #2
    \end{minipage}}%
  \vspace{0.1em}%
}

\lstdefinestyle{SAFoam}{
  backgroundcolor=\color{SAHypGray},
  basicstyle=\ttfamily\footnotesize,
  keywordstyle=\color{SAHypBlue}\bfseries,
  commentstyle=\color{SAHypSteel},
  numbers=left,
  numberstyle=\scriptsize\color{gray},
  numbersep=7pt,
  frame=single,
  rulecolor=\color{SAHypRule},
  breaklines=true,
  keepspaces=true,
  columns=fullflexible,
  showstringspaces=false,
  aboveskip=7pt,
  belowskip=7pt
}


\section{Summary and principal finding}
\label{sec:sa-summary}

This appendix audits the single-phase compressible Spalart--Allmaras (SA)
transport equation assembled by \SAOF{} v2412 (which is unchanged to the, as of writing, latest v2606).  It compares the native
finite-volume equation with the material-form equation cited by \SAOF{} from
Spalart and Allmaras (1994) and with its equivalent density-conservative form
given by Allmaras, Johnson, and Spalart (2012)
\cite{SpalartAllmaras1994,AJS2012}.

The issue is confined to the diffusion term.  Define
\[
  D_{\widetilde{\nu}}=\frac{\nu+\widetilde{\nu}}{\sigma}.
\]
Native \SAOF{} uses the density-weighted diffusive flux
\[
  \SAdiv\!\left(
    \rho D_{\widetilde{\nu}}\SAgrad\widetilde{\nu}
  \right),
\]
whereas multiplying the cited material equation by density requires
\[
  \rho\SAdiv\!\left(
    D_{\widetilde{\nu}}\SAgrad\widetilde{\nu}
  \right).
\]
These expressions are not equal for variable density.  The product rule gives
\[
  \SAdiv\!\left(
    \rho D_{\widetilde{\nu}}\SAgrad\widetilde{\nu}
  \right)
  =
  \rho\SAdiv\!\left(
    D_{\widetilde{\nu}}\SAgrad\widetilde{\nu}
  \right)
  +D_{\widetilde{\nu}}
   \SAgrad\rho\mathbin{\cdot}\SAgrad\widetilde{\nu}.
\]
The 2012 conservative equation cancels the final cross-gradient contribution
explicitly; native \SAOF{} does not.  Consequently, expanding the native
equation back into material form produces the additional term
\begin{equation}
  \boxed{
  \frac{\nu+\widetilde{\nu}}{\sigma}
  \SAgrad\ln\rho\mathbin{\cdot}\SAgrad\widetilde{\nu}}
  \label{eq:sa-summary-extra-term}
\end{equation}
relative to the cited 1994 equation.  Equivalently, relative to the 2012
conservative equation, the native physical right-hand side lacks
\begin{equation}
  \boxed{
  -\frac{\nu+\widetilde{\nu}}{\sigma}
  \SAgrad\rho\mathbin{\cdot}\SAgrad\widetilde{\nu}}.
  \label{eq:sa-summary-missing-term}
\end{equation}

\SAnote{SAHypLightBlue}{%
\begin{itemize}
  \item \textbf{Single identified mismatch.}  Native \SAOF{} places
    \(\rho\) inside the SA diffusive flux without the cross-gradient
    cancellation required to recover the cited material equation.
  \item \textbf{Two equivalent checks.}  The native material equation has
    the additional positive term in
    Eq.~\eqref{eq:sa-summary-extra-term}; equivalently, its conservative
    physical right-hand side lacks the negative term in
    Eq.~\eqref{eq:sa-summary-missing-term}.  These are two descriptions of
    the same continuum mismatch, not two independent inconsistencies.
  \item \textbf{Corrected implementation.}  The custom
    \SAcode{SpalartAllmarasDensityCorrected} model adds the required term with
    the coefficient and matrix sign needed to cancel the native
    density-product term.  It matches the 2012 diffusion coefficient when
    \SAcode{sigmaNut} is identified with the paper's \(\sigma\).
  \item \textbf{Scope.}  The mismatch vanishes for constant density or when
    \(\SAgrad\rho\cdot\SAgrad\widetilde{\nu}=0\).  The derivation concerns
    the continuum transport/diffusion sector; it does not establish complete
    closure identity, exact discrete equivalence, or improved predictions.
  \item \textbf{Independent NASA implementation warning.}  Note~4 of the
    NASA Turbulence Modeling Resource warns that some SA implementations
    ``incorrectly introduce the density inside the derivatives of the SA
    equation'' and states that this changes predictions for supersonic flows;
    it explicitly distinguishes the separately formulated Catris--Aupoix
    proposal as having a valid basis~\cite{NASATMRSA}.  The NASA note does not
    identify \SAOF{} by name.  Its relevance here follows from the independent
    reconstruction of Eq.~\eqref{eq:sa-native-conservative}, which establishes
    that the native implementation has precisely the warned-about derivative
    structure.
\end{itemize}
}

\subsection*{Outline}

\begin{itemize}
  \item Section~\ref{sec:sa-reference} identifies the literature equations,
    citation basis, and audit scope.
  \item Section~\ref{sec:sa-native-reconstruction} follows the actual
    \SAOF{} operators from code to the native \SAOF implemented conservative continuum equation.
  \item Section~\ref{sec:sa-comparison} directly compares the Allmaras 2012
    and native \SAOF{} conservative equations and identifies the missing term.
  \item Section~\ref{sec:sa-conclusion} summarizes the established algebraic
    result and interprets the native equation as a modified SA model.
  \item Section~\ref{sec:sa-correction} audits the custom correction,
    including the \SAcode{fvm::Su} sign and a dimensional check.
  \item Section~\ref{sec:sa-paper-derivation} retains the complete
    line-by-line derivation from the material equation to its
    density-conservative form.
\end{itemize}

\section{Reference equations and audit scope}
\label{sec:sa-reference}

The SA working variable is denoted by \(\widetilde{\nu}\), the molecular
kinematic viscosity by \(\nu\), density by \(\rho\), velocity by
\(\boldsymbol u\), and the model diffusion constant by \(\sigma\).  The
production, destruction, and trip contributions are grouped as \(P\), \(D\),
and \(T\), respectively.  The diffusion coefficient is left explicitly as
\((\nu+\widetilde{\nu})/\sigma\) so that every factor remains visible.

\subsection{The cited 1994 material formulation}

The audited \SAOF{} v2412 header cites the 1994 journal article by Spalart and
Allmaras as the standard-model reference.  That article explicitly defines
\(\mathrm D/\mathrm Dt\) as the Lagrangian or material derivative and gives the
full Version-Ia working-variable equation in Appendix Eq.~(A2), printed on
p.~20 (PDF page 16):
\begin{align}
  \SADt\widetilde{\nu}
  ={}&c_{b1}(1-f_{t2})\widetilde S\widetilde{\nu} \notag\\
  &+\frac{1}{\sigma}
   \left\{
     \SAdiv\!\left[(\nu+\widetilde{\nu})\SAgrad\widetilde{\nu}\right]
     +c_{b2}\left|\SAgrad\widetilde{\nu}\right|^2
   \right\} \notag\\
  &-\left(c_{w1}f_w-\frac{c_{b1}}{\kappa^2}f_{t2}\right)
    \left(\frac{\widetilde{\nu}}{d}\right)^2
  +f_{t1}(\Delta U)^2.
  \label{eq:sa-1994-material}
\end{align}
Production, destruction, and tripping differ in detail from the implementation
audited below, but the material transport and diffusion structure is the one
relevant to the density comparison.

The final paragraph of the 1994 appendix states that Eqs.~(A1)--(A6) had been
applied directly in compressible flow and reports only a weak influence of
density variations in the slightly supersonic boundary layers then examined.
Thus the cited journal article supports a density-free material equation for
compressible use.  It does not give a density-conservative finite-volume
recasting or prescribe where \(\rho\) should be placed within such a recasting.
The Turbulence Modeling Resource identifies this journal article as the
official publication arising from the 1992 conference paper and states that
the journal version takes precedence~\cite{SpalartAllmaras1994,NASATMRSA}.

\subsection{The Allmaras 2012 standard compressible form}
This resource explicitly states a conservative form of the SA model, which can be directly compared to the \SAOF{} implementation.
For this Allmaras et al. (2012) first write their transport equation Eq.~(2), printed on p.~2, as
\begin{equation}
  \SADt\widetilde{\nu}
  =P-D+T
  +\frac{1}{\sigma}
   \left\{
     \SAdiv\!\left[(\nu+\widetilde{\nu})\SAgrad\widetilde{\nu}\right]
     +c_{b2}\left|\SAgrad\widetilde{\nu}\right|^2
   \right\}.
  \label{eq:sa-paper-material}
\end{equation}
where production, wall destruction and trip terms are
\begin{equation*}
  P=c_{b1}(1-f_{t2})\widetilde{S}\widetilde{\nu},
  \qquad
  D=\left(c_{w1}f_w-\frac{c_{b1}}{\kappa^2}f_{t2}\right)
    \left[\frac{\widetilde{\nu}}{d}\right]^2,
  \qquad
  T=f_{t1}(\Delta u)^2.
\end{equation*}

This formulation is equivalent to the original 1994 model equation (see \ref{eq:sa-1994-material})

Their Section~2.2 explicitly states that this equation applies to both
incompressible and compressible flow and should be treated as the standard
compressible form.  They then give the equivalent conservation form as their
Eq.~(9), printed on p.~3:

\SAeqbox{SAHypLightBlue}{%
\textbf{density conservative SA model equation}
\begin{align}
  0={}&
  \frac{\partial(\rho\widetilde{\nu})}{\partial t}
  +\SAdiv(\rho\boldsymbol u\widetilde{\nu})
  -\rho(P-D+T) \notag\\
  &-\frac{1}{\sigma}
   \SAdiv\!\left[\rho(\nu+\widetilde{\nu})
   \SAgrad\widetilde{\nu}\right]
  -\frac{c_{b2}}{\sigma}\rho
   \left|\SAgrad\widetilde{\nu}\right|^2 \notag\\
  &+\frac{\nu+\widetilde{\nu}}{\sigma}
   \SAgrad\rho\mathbin{\cdot}\SAgrad\widetilde{\nu}.
  \label{eq:sa-paper-conservative}
\end{align}
}\par
The final term is positive because Eq.~\eqref{eq:sa-paper-conservative} is an
all-terms-left residual.  It is negative when moved to the physical
right-hand side.  Allmaras 2012 Eq.~(10), their shock-layer integral, independently
repeats the negative right-hand-side sign~\cite{AJS2012}.

Section~\ref{sec:sa-paper-derivation} contains the full mathematical derivation from the material equation Eq. (\ref{eq:sa-1994-material}) to the density conservative form Eq. (\ref{eq:sa-paper-conservative}) and proofs that no density external correction or similar is being applied, but that it is a clean mathematical derivation. Thus Eq. (\ref{eq:sa-paper-conservative}) should be understood as the faithful density conservative model equation of the original SA model.

\subsection{The 1992 Outlook on possible compressibility additions}

The earlier conference paper is retained here only for its historical Outlook.
That discussion mentions possible compressibility additions involving
\(\SAgrad\rho\cdot\SAgrad\widetilde{\nu}\), considers writing the transport
equation for \(\widetilde{\mu}\) instead of \(\widetilde{\nu}\), and notes
freedom to place density inside or outside derivatives.  It does not prescribe
a final coefficient, sign, or unique conservative equation, and therefore does
not define the density-weighted flux used by native \SAOF{} or establish its
equivalence to the standard material equation~\cite{SpalartAllmaras1992}.

\section{Reconstruction of the native OpenFOAM equation}
\label{sec:sa-native-reconstruction}

The governing equation is assembled in the shared SA base class.  The relevant
\SAOF{} v2412 code is shown first and subsequently the continuum equation is translated from it term-by-term to the underlying mathematical equation. Which can then be compared to the Allmaras 2012 published density conservative SA model equation. 
\Needspace{18\baselineskip}
\begin{lstlisting}[style=SAFoam,language=C++,firstnumber=468,
  caption={Native SA equation assembly in SpalartAllmarasBase.C.},
  label={lst:sa-native-assembly}]
fvm::ddt(alpha, rho, nuTilda_)
+ fvm::div(alphaRhoPhi, nuTilda_)
- fvm::laplacian(alpha*rho*DnuTildaEff(), nuTilda_)
- Cb2_/sigmaNut_*alpha()*rho()*magSqr(fvc::grad(nuTilda_)()())
==
  Cb1_*alpha()*rho()*Stilda()*nuTilda_()*(1 - ft2())
- fvm::Sp
  (
    (Cw1_*fw(Stilda, dTilda) - Cb1_/sqr(kappa_)*ft2())
   *alpha()*rho()*nuTilda_()/sqr(dTilda()),
    nuTilda_
  )
+ fvOptions(alpha, rho, nuTilda_)
\end{lstlisting}
This code is at
\SAcode{SpalartAllmarasBase.C:466--481}.  The coefficient used by the
Laplacian is defined at lines 383--389 as\par
\begin{equation}
  \SAcode{DnuTildaEff()}
  =\frac{\nu+\widetilde{\nu}}{\sigma}.
  \label{eq:sa-effective-diffusivity}
\end{equation}
Here and below, the code coefficient \SAcode{sigmaNut} is denoted by
\(\sigma\).  Its v2412 default is the configurable decimal
\SAcode{0.66666}; exact numerical agreement with the paper value
\(\sigma=2/3\) requires configuring the same value.  The custom correction
reuses \SAcode{DnuTildaEff} and therefore cancels the native density-product
term for either setting.

The thin RAS wrapper does not replace this equation; it supplies wall distance
and viscosity-related updates and inherits the base-class transport equation.

The reconstruction below is restricted to the bulk/interior transport
equation assembled in Listing~\ref{lst:sa-native-assembly}.  It precedes and
does not include the subsequent \SAcode{relax()},
\SAcode{fvOptions.constrain()}, \SAcode{boundaryManipulate()} (including any
boundary-row edits made by constrained patch fields), \SAcode{solve()}, the
post-solve \SAcode{bound()}, or
\SAcode{nuTilda\_.correctBoundaryConditions()} calls.  Those operations may
modify matrix coefficients, sources, boundary-adjacent rows, or the accepted
field, but in an audit it was not found that these operations correct the identified model deviation.

\subsection{Translate each code term into continuum notation}

For the registered single-phase compressible model, \(\alpha=1\), \(\rho\) is
the density field, and \SAcode{alphaRhoPhi} is the face mass flux corresponding
to \(\rho\boldsymbol u\cdot\boldsymbol S_f\).  The mappings below are for this
single-phase registration and interpret the displayed matrix as an
all-terms-left residual before reporting a physical right-hand-side sign.

\subsubsection{Transient and convection operators}

With \SAcode{fvm::ddt} and \SAcode{fvm::div} read as their canonical unbounded
continuum operators, the first two code terms,
\[
\begin{aligned}
  &\SAcode{fvm::ddt(alpha, rho, nuTilda\_)}
  \\
  &\quad+\SAcode{fvm::div(alphaRhoPhi, nuTilda\_)},
\end{aligned}
\]
represent
\begin{equation}
  \frac{\partial(\rho\widetilde{\nu})}{\partial t}
  +\SAdiv(\rho\boldsymbol u\widetilde{\nu}),
  \label{eq:sa-native-transport-operators}
\end{equation}
This is initially a conservative transport operator, not directly the
material derivative.  For differentiable fields, applying the product rule
gives
\begin{align}
  &\frac{\partial(\rho\widetilde{\nu})}{\partial t}
  +\SAdiv(\rho\boldsymbol u\widetilde{\nu}) \notag\\
  &\quad={}
  \rho\SADt\widetilde{\nu}
  +\widetilde{\nu}
  \left[
    \frac{\partial\rho}{\partial t}
    +\SAdiv(\rho\boldsymbol u)
  \right].
  \label{eq:sa-native-transport-product-rule}
\end{align}
Mass continuity,
\(\partial_t\rho+\SAdiv(\rho\boldsymbol u)=0\), removes the bracket.
Consequently,
\[
  \SAcode{ddt + div}
  \quad\longrightarrow\quad
  \rho\SADt\widetilde{\nu},
\]
and it becomes the unweighted material derivative only after the complete
equation is divided by \(\rho>0\).

This is a continuum/operator mapping, not a scheme-independent statement about
the assembled matrix.  The selected \SAcode{fvSchemes} may wrap the transient
or convection operator in a bounded form.  In particular, the bounded
three-field time scheme subtracts
\begin{center}
  \path{fvm::Sp(fvc::ddt(alpha,rho),vf)},
\end{center}
whereas the bounded convection scheme subtracts
\begin{center}
  \path{fvm::Sp(fvc::surfaceIntegrate(faceFlux),vf)}.
\end{center}
These terms change how a discrete mass-continuity residual is apportioned
between transport and source; they do not contain the diffusion cross-gradient
\((\nu+\widetilde{\nu})\SAgrad\rho\cdot\SAgrad\widetilde{\nu}/\sigma\).

\subsubsection{Laplacian operator}

Using Eq.~\eqref{eq:sa-effective-diffusivity}, the matrix-left term
\begin{equation*}
  -\SAcode{fvm::laplacian}
  \left(
    \rho\frac{\nu+\widetilde{\nu}}{\sigma},
    \widetilde{\nu}
  \right)
\end{equation*}
represents
\begin{equation}
  -\frac{1}{\sigma}
  \SAdiv\!\left[
    \rho(\nu+\widetilde{\nu})\SAgrad\widetilde{\nu}
  \right].
  \label{eq:sa-native-laplacian}
\end{equation}
Density is inside the coefficient passed to the Laplacian.  The code therefore
discretizes a density-weighted conservative diffusive flux; it does not
discretize \(\rho\) multiplied by a density-free Laplacian.

\subsubsection{Explicit \texorpdfstring{\(c_{b2}\)}{cb2} term}

The code term
\[
  -\SAcode{Cb2\_/sigmaNut\_*rho()*magSqr(fvc::grad(nuTilda\_)()())}
\]
maps to
\begin{equation}
  -\frac{c_{b2}}{\sigma}\rho
  \left|\SAgrad\widetilde{\nu}\right|^2.
  \label{eq:sa-native-cb2}
\end{equation}
The \(c_{b2}\) term is part of the modelled nonlinear diffusion in both
papers.  It is distinct from the density product-rule term and must not be
interpreted as the missing cancellation.

\subsubsection{Production and destruction terms}

The first term on the code right-hand side represents the paper's production
term:
\[
  \SAcode{Cb1\_*rho()*Stilda()*nuTilda\_()*(1-ft2())}
  \quad\longrightarrow\quad
  \rho P,
\]
where
\begin{equation}
  P=c_{b1}(1-f_{t2})\widetilde S\widetilde{\nu}.
  \label{eq:sa-native-production}
\end{equation}
The implicit \SAcode{fvm::Sp} term represents destruction.  Because it is
preceded by a minus sign on the code right-hand side, its continuum meaning is
\(-\rho D\), where
\begin{equation}
  D=
  \left(
    c_{w1}f_w-\frac{c_{b1}}{\kappa^2}f_{t2}
  \right)
  \left(\frac{\widetilde{\nu}}{d}\right)^2.
  \label{eq:sa-native-destruction}
\end{equation}
Thus the modelled right-hand side in the native equation is \(\rho(P-D)\).
Native \SAOF{} has no trip term \(T\); comparison with the full paper equation
therefore uses \(T=0\).  The final
\SAcode{fvOptions(alpha, rho, nuTilda\_)} entry permits additional case-defined
sources.  It is not part of the SA paper equation and is omitted from the
reconstructed native equation.

For the density audit, \(P\) and \(D\) are held symbolically fixed so that the
transport and diffusion sectors can be compared directly.  This does not
assert complete closure identity: native \SAOF{} uses the
\(\widetilde S\) limiter
\(\max[\Omega+f_{v2}\widetilde{\nu}/(\kappa d)^2,C_s\Omega]\), which is not
identical to the Allmaras 2012 baseline and proposed negative-\(\widetilde S\)
continuation.

\subsection{Assemble the mathematical equation}

Combining the preceding mappings gives the all-terms-left continuum meaning of
the native finite-volume assembly:
\begin{align}
  0={}&
  \frac{\partial(\rho\widetilde{\nu})}{\partial t}
  +\SAdiv(\rho\boldsymbol u\widetilde{\nu})
  -\rho(P-D) \notag\\
  &-\frac{1}{\sigma}
    \SAdiv\!\left[
      \rho(\nu+\widetilde{\nu})\SAgrad\widetilde{\nu}
    \right]
  -\frac{c_{b2}}{\sigma}\rho
    \left|\SAgrad\widetilde{\nu}\right|^2.
  \label{eq:sa-native-residual}
\end{align}
Equivalently, moving the modelled and diffusion terms to the physical
right-hand side gives:

\SAeqbox{SAHypLightBlue}{%
\textbf{Native OpenFOAM continuum equation}
\begin{align}
  \frac{\partial(\rho\widetilde{\nu})}{\partial t}
  +\SAdiv(\rho\boldsymbol u\widetilde{\nu})
  ={}&\rho(P-D) \notag\\
  &+\frac{1}{\sigma}
    \SAdiv\!\left[
      \rho(\nu+\widetilde{\nu})\SAgrad\widetilde{\nu}
    \right]
  +\frac{c_{b2}}{\sigma}\rho
    \left|\SAgrad\widetilde{\nu}\right|^2.
  \label{eq:sa-native-conservative}
\end{align}
}\par
Equation~\eqref{eq:sa-native-conservative} is not copied from either paper; it
is the continuum equation represented by the audited native finite-volume
assembly, expressed in the papers' notation with \(\alpha=1\),
\SAcode{fvOptions} suppressed, and \(T=0\).  For the incompressible
registration, both \(\alpha\) and \(\rho\) are compile-time one fields, so the
density issue vanishes exactly.

\Needspace{0.47\textheight}
\section{Direct comparison and identified error}
\label{sec:sa-comparison}

Equation~\eqref{eq:sa-paper-conservative} is the Allmaras 2012 conservative
equation, and Eq.~\eqref{eq:sa-native-conservative} is the equation
reconstructed term by term from native \SAOF{}.  Set \(T=0\), as required by
the native implementation, and hold \(P\), \(D\), and \(\sigma\) fixed.
After moving the modelled and diffusion terms of
Eq.~\eqref{eq:sa-paper-conservative} to the physical right-hand side, the
complete comparison is
\begin{equation}
  \left.\mathrm{RHS}\right|_{\text{Allmaras Eq.~\eqref{eq:sa-paper-conservative}},\,T=0}
  =
  \left.\mathrm{RHS}\right|_{\text{OpenFOAM Eq.~\eqref{eq:sa-native-conservative}}}
  -\frac{\nu+\widetilde{\nu}}{\sigma}
   \SAgrad\rho\mathbin{\cdot}\SAgrad\widetilde{\nu}.
  \label{eq:sa-direct-comparison}
\end{equation}
Every other displayed transport, production, destruction, and diffusion term
is identical under this comparison.  Equation~\eqref{eq:sa-native-conservative}
therefore lacks the final negative density-gradient term.

That term is the cancellation required by the product rule
\[
  \SAdiv\!\left[
    \rho(\nu+\widetilde{\nu})\SAgrad\widetilde{\nu}
  \right]
  =
  \rho\SAdiv\!\left[
    (\nu+\widetilde{\nu})\SAgrad\widetilde{\nu}
  \right]
  +(\nu+\widetilde{\nu})
   \SAgrad\rho\mathbin{\cdot}\SAgrad\widetilde{\nu}.
\]
Without the cancellation, rewriting the native equation in material form
produces the corresponding positive cross-gradient term.

\SAeqbox{SAHypLightRed}{%
\textbf{Identified error in the \SAOF implemented equation.}

Relative to Eq.~\eqref{eq:sa-paper-conservative}, the physical right-hand side
of Eq.~\eqref{eq:sa-native-conservative} lacks
\[
  -\frac{\nu+\widetilde{\nu}}{\sigma}
    \SAgrad\rho\mathbin{\cdot}\SAgrad\widetilde{\nu}.
\]
Equivalently, the native material equation contains the additional term
\[
  +\frac{\nu+\widetilde{\nu}}{\sigma}
    \SAgrad\ln\rho\mathbin{\cdot}\SAgrad\widetilde{\nu}.
\]
}

The equations agree when density is constant,
\(\SAgrad\widetilde{\nu}=0\), or the density and working-variable gradients
are orthogonal.  A nonzero density gradient by itself is not sufficient; the
cross-gradient alignment determines the contribution.

\section{Conclusion}
\label{sec:sa-conclusion}

\SAnote{SAHypLightGreen}{%
Native single-phase compressible \SAOF{} v2412 uses the density-weighted
diffusive flux
\[
  \frac{1}{\sigma}\SAdiv\!\left[
    \rho(\nu+\widetilde{\nu})\SAgrad\widetilde{\nu}
  \right]
\]
without the explicit conservative cancellation needed to recover
\[
  \frac{\rho}{\sigma}\SAdiv\!\left[
    (\nu+\widetilde{\nu})\SAgrad\widetilde{\nu}
  \right].
\]
Its material form therefore contains the additional positive term in
Eq.~\eqref{eq:sa-summary-extra-term}.  The custom
\SAcode{SpalartAllmarasDensityCorrected} source adds the corresponding positive
term to the matrix residual, which is negative on the physical right-hand side
and cancels the native product-rule contribution.  With the same \(\sigma\),
the resulting diffusion structure agrees with Allmaras 2012 Eq.~(9).  This
    cancellation is exact at the continuum level under the assumptions stated
    during the reconstruction in Section~\ref{sec:sa-native-reconstruction};
    it is not a claim that all nondiffusive closure choices are identical.
}

Equation~\eqref{eq:sa-native-conservative} nevertheless defines a
mathematically consistent conservative scalar-transport model if the diffusive
flux
\begin{equation}
  \boldsymbol J_{\widetilde{\nu}}
  =-\rho D_{\widetilde{\nu}}\SAgrad\widetilde{\nu},
  \qquad
  D_{\widetilde{\nu}}\equiv\frac{\nu+\widetilde{\nu}}{\sigma},
  \label{eq:sa-native-fickian-flux}
\end{equation}
is chosen intentionally.  That choice, however, generates the additional
density-gradient term in the corresponding material equation.  Without the
cancellation prescribed by Allmaras et al. (2012), the native equation is
therefore a modified SA model rather than the intended compressible SA
formulation.

The defensible conclusion is thus a \emph{confirmed
transport/diffusion-sector deviation from both the 1994 material formulation
cited by \SAOF{} and the Allmaras 2012 conservative standard}.  This audit
cannot determine whether the deviation was intentional or arose from an
unintended derivation error.  The 1992 Outlook shows that alternative
compressibility additions and density placements were contemplated, but it
does not define the native density-weighted flux or establish that flux as the
intended SA model~\cite{SpalartAllmaras1994,SpalartAllmaras1992,AJS2012}.  The
native \SAOF{} equation and the cited standard therefore represent different
models in compressible flow, with the latter and herein corrected, representing a version in-line with the understanding of Spalart--Allmaras and stated as the standard form for compressible flow \cite{AJS2012}.

\clearpage
\section{Implementation of the corrected model}
\label{sec:sa-correction}

The custom model first assembles the native equation and then adds the source
shown below.  The relevant code is in
\SAcode{SpalartAllmarasDensityCorrected.C:54--66,136--175}.
\begin{lstlisting}[style=SAFoam,language=C++,firstnumber=54,
  caption={Density-gradient source in SpalartAllmarasDensityCorrected.},
  label={lst:sa-custom-correction}]
const tmp<volVectorField> tgradRho = fvc::grad(rho);

// The stock conservative diffusion contributes
//     div(rho*DnuTildaEff*grad(nuTilda))
// on the physical RHS.  Adding this explicit source to the matrix LHS
// subtracts DnuTildaEff*grad(rho).grad(nuTilda) from that RHS, giving
// exactly rho*div(DnuTildaEff*grad(nuTilda)).
nuTildaEqn += fvm::Su
(
    DnuTildaEff()
   *(tgradRho()() & gradNuTilda()),
    this->nuTilda_
);
\end{lstlisting}

The \SAcode{fvm::Su} convention resolves the apparent sign ambiguity.  The
explicit calls to the field-internal accessors make the expression in
Listing~\ref{lst:sa-custom-correction} a temporary
\SAcode{DimensionedField}.  The overload in \SAcode{fvmSup.C:92--103}
forwards that temporary to the field implementation at lines 72--89, where the
supplied field is stored as \SAcode{matrix.source() -= V*source}.  Adding the
source therefore adds
\begin{equation}
  +\frac{\nu+\widetilde{\nu}}{\sigma}
  \SAgrad\rho\mathbin{\cdot}\SAgrad\widetilde{\nu}
  \label{eq:sa-custom-residual-source}
\end{equation}
to the all-terms-left matrix residual, or the negative of this expression to
the physical right-hand side.  It has the form and sign of the final term in
Eq.~\eqref{eq:sa-paper-conservative}, is exactly the product-rule cancellation
identified in Section~\ref{sec:sa-comparison} for the native
\SAcode{DnuTildaEff}, and has the Allmaras 2012 numerical coefficient when
\SAcode{sigmaNut} is set to \(2/3\).

The dimensions are consistent:
\begin{equation}
  \left[\frac{\nu+\widetilde{\nu}}{\sigma}\right]=L^2T^{-1},
  \qquad
  [\SAgrad\rho]=ML^{-4},
  \qquad
  [\SAgrad\widetilde{\nu}]=LT^{-1},
  \label{eq:sa-custom-source-factors-dimensions}
\end{equation}
and hence
\begin{equation}
  \left[
    \frac{\nu+\widetilde{\nu}}{\sigma}
    \SAgrad\rho\mathbin{\cdot}\SAgrad\widetilde{\nu}
  \right]
  =ML^{-1}T^{-2}
  =\left[\frac{\partial(\rho\widetilde{\nu})}{\partial t}\right].
  \label{eq:sa-custom-source-dimensions}
\end{equation}

The custom header provides a no-op overload for
\SAcode{geometricOneField} density, so the incompressible registration remains
the native SA equation.  The library is registered only in the single-phase
compressible and incompressible turbulence tables.  Its omission of a
phase-fraction-gradient term is consistent with that scope; applying the same
source to a variable-\(\alpha\) multiphase equation would require a new
derivation.

Finally, the modelled term
\begin{equation}
  \frac{c_{b2}}{\sigma}
  \left|\SAgrad\widetilde{\nu}\right|^2
\end{equation}
is not the density correction.  Expanding
\(\SAdiv[(\nu+\widetilde{\nu})\SAgrad\widetilde{\nu}]\) already retains the
derivatives of \(\nu+\widetilde{\nu}\), while \(c_{b2}\) is a separate model
coefficient.  The custom cancellation neither removes nor duplicates it.

\clearpage
\section{Allmaras 2012 material form to conservative form}
\label{sec:sa-paper-derivation}

This section retains the complete conversion from Allmaras 2012 Eq.~(2) to
Eq.~(9).  It makes explicit which terms are supplied by mass conservation and
which term is created by rewriting the diffusion operator.

The Allmaras 2012 material equation is Eq.~\eqref{eq:sa-paper-material}:
\begin{equation}
  \SADt\widetilde{\nu}
  =P-D+T
  +\frac{1}{\sigma}
   \left[
     \SAdiv\!\left((\nu+\widetilde{\nu})\SAgrad\widetilde{\nu}\right)
     +c_{b2}\left|\SAgrad\widetilde{\nu}\right|^2
   \right].
  \label{eq:sa-deriv-paper-start}
\end{equation}

\subsection{Put the material equation in zero-residual form}

Move every right-hand-side term to the left and expand the material derivative:
\begin{equation}
\begin{aligned}
  \{\mathrm{S\!A}\}={}&
  \frac{\partial\widetilde{\nu}}{\partial t}
  +\boldsymbol u\mathbin{\cdot}\SAgrad\widetilde{\nu}
  -(P-D+T) \\
  &-\frac{1}{\sigma}
   \SAdiv\!\left[(\nu+\widetilde{\nu})\SAgrad\widetilde{\nu}\right]
  -\frac{c_{b2}}{\sigma}
   \left|\SAgrad\widetilde{\nu}\right|^2
  =0.
\end{aligned}
\label{eq:sa-deriv-sa-residual}
\end{equation}
The braces label the complete residual; they introduce no new model quantity.

\subsection{Add the zero mass residual}

Mass conservation in zero-residual form is
\begin{equation}
  \{\mathrm{mass}\}
  =\frac{\partial\rho}{\partial t}
  +\SAdiv(\rho\boldsymbol u)
  =0.
  \label{eq:sa-deriv-mass-residual}
\end{equation}
Allmaras et al. (2012) combine the two residuals as
\begin{equation}
  0=\rho\{\mathrm{S\!A}\}
  +\widetilde{\nu}\{\mathrm{mass}\}.
  \label{eq:sa-deriv-combined-residuals}
\end{equation}
Because both residuals are zero, multiplying and adding them does not change
the solution.

\Needspace{0.38\textheight}
\subsection{Multiply out the residuals}

Multiplying Eq.~\eqref{eq:sa-deriv-sa-residual} by \(\rho\) gives
\begin{align}
  \rho\{\mathrm{S\!A}\}={}&
  \rho\frac{\partial\widetilde{\nu}}{\partial t}
  +\rho\boldsymbol u\mathbin{\cdot}\SAgrad\widetilde{\nu}
  -\rho(P-D+T) \notag\\
  &-\frac{\rho}{\sigma}
   \SAdiv\!\left[(\nu+\widetilde{\nu})\SAgrad\widetilde{\nu}\right]
  -\frac{c_{b2}}{\sigma}\rho
   \left|\SAgrad\widetilde{\nu}\right|^2.
  \label{eq:sa-deriv-rho-sa}
\end{align}
Multiplying Eq.~\eqref{eq:sa-deriv-mass-residual} by
\(\widetilde{\nu}\) gives
\begin{equation}
  \widetilde{\nu}\{\mathrm{mass}\}
  =\widetilde{\nu}\frac{\partial\rho}{\partial t}
  +\widetilde{\nu}\SAdiv(\rho\boldsymbol u).
  \label{eq:sa-deriv-nutilda-mass}
\end{equation}
Adding Eqs.~\eqref{eq:sa-deriv-rho-sa} and
\eqref{eq:sa-deriv-nutilda-mass} yields
\begin{align}
  0={}&
  \rho\frac{\partial\widetilde{\nu}}{\partial t}
  +\widetilde{\nu}\frac{\partial\rho}{\partial t} \notag\\
  &+\rho\boldsymbol u\mathbin{\cdot}\SAgrad\widetilde{\nu}
  +\widetilde{\nu}\SAdiv(\rho\boldsymbol u)
  -\rho(P-D+T) \notag\\
  &-\frac{\rho}{\sigma}
   \SAdiv\!\left[(\nu+\widetilde{\nu})\SAgrad\widetilde{\nu}\right]
  -\frac{c_{b2}}{\sigma}\rho
   \left|\SAgrad\widetilde{\nu}\right|^2.
  \label{eq:sa-deriv-expanded-sum}
\end{align}

\subsection{Combine transient and convection products}

For the transient terms,
\begin{equation}
  \frac{\partial(\rho\widetilde{\nu})}{\partial t}
  =\rho\frac{\partial\widetilde{\nu}}{\partial t}
  +\widetilde{\nu}\frac{\partial\rho}{\partial t}.
  \label{eq:sa-deriv-transient-product}
\end{equation}
For a scalar \(\widetilde{\nu}\) and vector \(\rho\boldsymbol u\),
\begin{equation}
  \SAdiv(\rho\boldsymbol u\widetilde{\nu})
  =\rho\boldsymbol u\mathbin{\cdot}\SAgrad\widetilde{\nu}
  +\widetilde{\nu}\SAdiv(\rho\boldsymbol u).
  \label{eq:sa-deriv-convection-product}
\end{equation}
Substitution into Eq.~\eqref{eq:sa-deriv-expanded-sum} gives
\begin{align}
  0={}&
  \frac{\partial(\rho\widetilde{\nu})}{\partial t}
  +\SAdiv(\rho\boldsymbol u\widetilde{\nu})
  -\rho(P-D+T) \notag\\
  &-\frac{\rho}{\sigma}
   \SAdiv\!\left[(\nu+\widetilde{\nu})\SAgrad\widetilde{\nu}\right]
  -\frac{c_{b2}}{\sigma}\rho
   \left|\SAgrad\widetilde{\nu}\right|^2.
  \label{eq:sa-deriv-before-diffusion-rewrite}
\end{align}
Only the diffusion term remains to be rewritten.

\subsection{Rewrite the diffusion term}

Apply the vector product rule:
\begin{align}
  &\SAdiv\!\left[
    \rho(\nu+\widetilde{\nu})\SAgrad\widetilde{\nu}
  \right] \notag\\
  &\quad={}
  \rho\SAdiv\!\left[
    (\nu+\widetilde{\nu})\SAgrad\widetilde{\nu}
  \right]
  +(\nu+\widetilde{\nu})
   \SAgrad\rho\mathbin{\cdot}\SAgrad\widetilde{\nu}.
  \label{eq:sa-deriv-diffusion-product}
\end{align}
Solve for the density-multiplied diffusion that appears in
Eq.~\eqref{eq:sa-deriv-before-diffusion-rewrite}:
\begin{align}
  &\rho\SAdiv\!\left[
    (\nu+\widetilde{\nu})\SAgrad\widetilde{\nu}
  \right] \notag\\
  &\quad={}
  \SAdiv\!\left[
    \rho(\nu+\widetilde{\nu})\SAgrad\widetilde{\nu}
  \right]
  -(\nu+\widetilde{\nu})
   \SAgrad\rho\mathbin{\cdot}\SAgrad\widetilde{\nu}.
  \label{eq:sa-deriv-diffusion-rearranged}
\end{align}
The term in Eq.~\eqref{eq:sa-deriv-before-diffusion-rewrite} is negative, so
\begin{align}
  &-\frac{\rho}{\sigma}\SAdiv\!\left[
    (\nu+\widetilde{\nu})\SAgrad\widetilde{\nu}
  \right] \notag\\
  &\quad={}
  -\frac{1}{\sigma}\SAdiv\!\left[
    \rho(\nu+\widetilde{\nu})\SAgrad\widetilde{\nu}
  \right]
  +\frac{\nu+\widetilde{\nu}}{\sigma}
   \SAgrad\rho\mathbin{\cdot}\SAgrad\widetilde{\nu}.
  \label{eq:sa-deriv-diffusion-residual-sign}
\end{align}
The positive final sign follows from two negative signs: the one already in
front of the residual diffusion and the one introduced by rearranging the
product rule.

\subsection{Recover Allmaras 2012 Eq.~(9)}

Substituting Eq.~\eqref{eq:sa-deriv-diffusion-residual-sign} into
Eq.~\eqref{eq:sa-deriv-before-diffusion-rewrite} gives
\begin{align}
  0={}&
  \frac{\partial(\rho\widetilde{\nu})}{\partial t}
  +\SAdiv(\rho\boldsymbol u\widetilde{\nu})
  -\rho(P-D+T) \notag\\
  &-\frac{1}{\sigma}\SAdiv\!\left[
    \rho(\nu+\widetilde{\nu})\SAgrad\widetilde{\nu}
  \right]
  -\frac{c_{b2}}{\sigma}\rho
   \left|\SAgrad\widetilde{\nu}\right|^2 \notag\\
  &+\frac{\nu+\widetilde{\nu}}{\sigma}
   \SAgrad\rho\mathbin{\cdot}\SAgrad\widetilde{\nu}.
  \label{eq:sa-deriv-paper-final}
\end{align}
This is Allmaras 2012 Eq.~(9).  The mass residual creates conservative storage and
convection.  The diffusion product rule, not continuity, creates the final
density-gradient term.

\clearpage
\endgroup
